Whether professional basketball careers fall into discrete trajectory types or vary continuously matters for how longitudinal performance data should be modelled. We examine it for \nPlayers{} NBA players who debuted between 1979--80 and 2010--11, using season-level Player Efficiency Rating (PER) sequences and a protocol built around reproducibility. A Transformer autoencoder compresses variable-length careers almost losslessly (held-out RMSE \rmseHeld{} PER points against \rmseLinear{} for a player-specific linear trend), but its latent geometry depends on the training seed: a single UMAP+HDBSCAN run never reaches ARI 0.8 across seeds, and a five-seed consensus is reproducible only at three coarse groups. Under an identical consensus protocol, nine summary statistics are far more reproducible (consensus Silhouette \summSilEight{} versus \aeAllSilEight{} at $k=8$). Reproducibility is not evidence of types, however: representations disagree (ARI $\le$ \xMax), the gap statistic finds no structure beyond a correlated cloud, and unimodality is not rejected. Career productivity paths form a continuum summarised by three coordinates, level (\pcOneVar\% of variance), rise versus decline (\pcTwoVar\%) and length; we report an eight-type segmentation of it, associated with All-Star selection, Win Shares and draft status but not position, and withdraw an earlier claim of nineteen discrete types. Code and data are released.